\documentclass[10pt]{article}
\usepackage[utf8]{inputenc}

\usepackage[margin=1in]{geometry}
\usepackage{amsmath}
\usepackage{amssymb}
\usepackage{amsthm}
\usepackage{mathtools}
\usepackage{bm}
\usepackage{tikz}

\usepackage{color}
\usepackage{colortbl}
\definecolor{deepblue}{rgb}{0,0,0.5}
\definecolor{deepred}{rgb}{0.6,0,0}
\definecolor{deepgreen}{rgb}{0,0.5,0}
\definecolor{gray}{rgb}{0.7,0.7,0.7}

\usepackage{hyperref}
\hypersetup{
  colorlinks   = true, %Colours links instead of ugly boxes
  urlcolor     = black, %Colour for external hyperlinks
  linkcolor    = blue, %Colour of internal links
  citecolor    = blue  %Colour of citations
}

%%%%%%%%%%%%%%%%%%%%%%%%%%%%%%%%%%%%%%%%%%%%%%%%%%%%%%%%%%%%%%%%%%%%%%%%%%%%%%%%

\theoremstyle{definition}
\newtheorem{problem}{Problem}
\newtheorem{example}[problem]{Example}
\newtheorem{lemma}[problem]{Lemma}
\newtheorem{note}[problem]{Note}
\newtheorem{defn}[problem]{Definition}
\newcommand{\E}{\mathbb E}
\newcommand{\R}{\mathbb R}
\DeclareMathOperator{\nnz}{nnz}
\DeclareMathOperator{\determinant}{det}
\DeclareMathOperator{\Var}{Var}
\DeclareMathOperator{\rank}{rank}
\DeclareMathOperator*{\argmin}{arg\,min}
\DeclareMathOperator*{\argmax}{arg\,max}

\newcommand{\trans}[1]{{#1}^{T}}
\newcommand{\loss}{\ell}
\newcommand{\w}{\mathbf w}
\newcommand{\x}{\mathbf x}
\newcommand{\y}{\mathbf y}
\newcommand{\lone}[1]{{\lVert {#1} \rVert}_1}
\newcommand{\ltwo}[1]{{\lVert {#1} \rVert}_2}
\newcommand{\lp}[1]{{\lVert {#1} \rVert}_p}
\newcommand{\linf}[1]{{\lVert {#1} \rVert}_\infty}
\newcommand{\lF}[1]{{\lVert {#1} \rVert}_F}

\usepackage{listings}

% Default fixed font does not support bold face
\DeclareFixedFont{\ttb}{T1}{txtt}{bx}{n}{12} % for bold
\DeclareFixedFont{\ttm}{T1}{txtt}{m}{n}{12}  % for normal

% Python style for highlighting
\newcommand\pythonstyle{\lstset{
language=Python,
basicstyle=\ttm,
otherkeywords={self},             % Add keywords here
keywordstyle=\ttb\color{deepblue},
emph={MyClass,__init__},          % Custom highlighting
emphstyle=\ttb\color{deepred},    % Custom highlighting style
stringstyle=\color{deepgreen},
frame=tb,                         % Any extra options here
showstringspaces=false,           % 
stepnumber=1,
numbers=left
}}

\lstnewenvironment{python}[1][]
{
    \pythonstyle
    \lstset{#1}
}
{}

%%%%%%%%%%%%%%%%%%%%%%%%%%%%%%%%%%%%%%%%%%%%%%%%%%%%%%%%%%%%%%%%%%%%%%%%%%%%%%%%

\begin{document}

\begin{center}
    {
\Large
    Math: Asymptotic (Big-O/$\Theta$/$\Omega$) Notation in a Single Variable
}
\end{center}
\begin{note}
    We will use asymptotic notation to measure properties of our programs like: runtime, memory usage, and API cost.
    In future classes, you will also be using asymptotic notation for measuring statistical properties of programs as well like how accurate a machine learning algorithm is based on the number of data points trained on.
    We will therefore use a slightly technical definition of asymptotic notation that can be used in all of these scenarios.
    But we will focus on intuition and not get into weird mathematical technicalities that don't come up in practice.
\end{note}

\begin{defn}
    Let $f,g$ be functions from $\R^+\to\R^+$.
    Then,
    \begin{enumerate}
        \item If $\displaystyle\lim_{x\to\infty} \frac{f(x)}{g(x)} < \infty$, then we say $f = O(g)$.
        \item If $\displaystyle\lim_{x\to\infty} \frac{f(x)}{g(x)} > 0$, then we say $f = \Omega(g)$.
        \item We say that $f = \Theta(g)$ if both $f=O(g)$ and $f=\Omega(g)$.
    \end{enumerate}
    Intuitively, you should think of $O$ as $\le$, $\Omega$ as $\ge$, and $\Theta$ as $=$.
\end{defn}

\begin{example}~
    \begin{enumerate}
        \item 
            $f(x) = x$
            
            $g(x) = x^2$
            \vspace{0.6in}
        \item
            $f(x) = x^2$
            
            $g(x) = x$
            \vspace{0.6in}
        \item
            $f(x) = x^2 + 2x + 5$
            
            $g(x) = x$
            \vspace{0.6in}
        \item
            $f(x) = x^2 + 2x + 5$
            
            $g(x) = x^2$
            \vspace{0.6in}
        \item
            $f(x) = x^2 + 2x + 5$
            
            $g(x) = x^3$
    \end{enumerate}
\end{example}

\newpage

\begin{note}
    The main idea of asymptotic notation is that we can ``discard'' the ``small'' terms and only focus on the ``large'' terms.
    This makes working with functions much simpler.
    We lose ``accuracy'' for small numbers,
    but in practice the behavior of large numbers is much more important.
\end{note}
\begin{center}
\begin{tikzpicture}
\draw[help lines, color=gray!30, dashed] (-0.5,-0.5) grid (14.9,9.9);
\draw[-,ultra thick] (-0.5,0)--(15,0) node[right]{};
\draw[-,ultra thick] (0,-0.5)--(0,10) node[above]{};
\end{tikzpicture}
\end{center}

%\newpage
%\begin{example}
    %Complete each equation below by adding the symbol $O$ if $f=O(g)$, $\Omega$ if $f=\Omega(g)$, or $\Theta$ if $f=\Theta(g)$.  
    %The first row is completed for you as an example.
%
%{\renewcommand{\arraystretch}{4.4}
%\begin{tabular}{c c c c c c}
    %& $f(n)$ &~\hspace{0.5in}~$ $~\hspace{0.5in}~& $g(n)$ &\\
    %\hline
    %& $1$ & ~\hspace{0.5in}~$=$~\hspace{0.5in}~  & $O(n)$ &  &\\
    %\arrayrulecolor{gray}\hline
    %& $3 n\log n$ & ~\hspace{0.5in}~$=$~\hspace{0.5in}~  & $n^2$ &  &\\
    %\arrayrulecolor{gray}\hline
    %& $1$ & ~\hspace{0.5in}~$=$~\hspace{0.5in}~  & $1/n$ &  &\\
    %\arrayrulecolor{gray}\hline
    %& $\log_2 n$ & ~\hspace{0.5in}~$=$~\hspace{0.5in}~  & $\log_3 n$ &  &\\
    %\arrayrulecolor{gray}\hline
    %& $\log n$ & ~\hspace{0.5in}~$=$~\hspace{0.5in}~  & $\frac {1} {\log n}$ &  &\\
    %\arrayrulecolor{gray}\hline
    %& $5\cdot10^{30}$ & ~\hspace{0.5in}~$=$~\hspace{0.5in}~  & $\log n$ &  &\\
    %\arrayrulecolor{gray}\hline
    %& $\log n$ & ~\hspace{0.5in}~$=$~\hspace{0.5in}~  & $\log (n^2)$ &  &\\
    %\arrayrulecolor{gray}\hline
    %& $2^n$ & ~\hspace{0.5in}~$=$~\hspace{0.5in}~  & $3^n$ &  &\\
    %\arrayrulecolor{gray}\hline
    %& $\frac 1 n$ & ~\hspace{0.5in}~$=$~\hspace{0.5in}~  & $\sqrt{\frac 1 n}$ &  &\\
    %\arrayrulecolor{gray}\hline
    %& $\log n$ & ~\hspace{0.5in}~$=$~\hspace{0.5in}~  & $(\log n)^2$ &  &\\
    %\arrayrulecolor{gray}\hline
%\end{tabular}
%}
%\end{example}

\newpage
%\begin{defn}
    %An expression $E$ is \emph{simplified} if
    %\begin{enumerate}
        %\item $E$ contains no additive terms (i.e.\ $E$ is a single product of atoms), and
        %\item $E$ contains no constant factors other than the number $1$ itself.
    %\end{enumerate}
    %Equivalently: $E$ is as short as possible among all expressions $E'$ with $E=\Theta(E')$.
%\end{defn}

\begin{example}
It is often useful to simplify asymptotic expressions.
In this example, we will work out how to simplify
$$
    O\bigg((2n+1)^3\bigg)
$$
The procedure is first expand the expression into a sum of terms:
%\begin{enumerate}
%
    %\item Expand $E$ algebraically into a sum of terms.
%
        %$$
        %(2n+1)^3 = 8n^3 + 12n^2 + 6n + 1
        %$$
%
    %\item Replace leading constants with 1.
        %$$
        %\to n^3 + n^2 + n + 1
        %$$
%
    %\item Discard every term that is $O$ of another term.
%
        %$$
        %\to n^3
        %$$
%\end{enumerate}
\begin{align*}
  (2n+1)^3 &= 8n^3 + 12n^2 + 6n + 1
  \intertext{Replace leading constants with 1:}
           &\sim n^3 + n^2 + n + 1
  \intertext{Discard every term that is $O$ of another term:}
           &\sim n^3
\end{align*}
So our initial expression simplifies to $O(n^3)$.
This procedure works for $O$, $\Theta$, and $\Omega$.
\end{example}

\begin{example}
    Simplify the following expressions:

\begin{enumerate}
    \item $O\bigg(8n^3 + n^2 \bigg)$
    \vspace{0.8in}
    \item $\Omega\bigg((7n)^3 + 5n^2\log n + \log n \bigg)$
    \vspace{0.8in}
    \item $\Theta\bigg(100000000000 \bigg)$
    \vspace{0.8in}
    \item $O\bigg(\log n + 100000000000 \bigg)$
    \vspace{0.8in}
\item $O\bigg(\frac 1 n + \frac 1 {n^2}\bigg)$
    \vspace{0.8in}
%\item $O\bigg(ab^2 + 3a^2b + ab + 10b\bigg)$
    %\vspace{1.5in}
%\item $O\bigg(a + b + c + ab + bc + abc\bigg)$
    %\vspace{1.5in}
%\item $O\bigg(\frac 1 n + \frac 1 {nm} + \frac 1 m \bigg)$
    %\vspace{1.5in}
%\item $\Omega\bigg((3.45n + n)(\log n^2)\bigg)$
    %\vspace{1.5in}
%\item $\Theta\bigg(n(1 + \log n) + n^{3.2} + \log 2^n\bigg)$
    %\vspace{1.5in}
\end{enumerate}
\end{example}

\vspace{0.25in}
\noindent

%%%%%%%%%%%%%%%%%%%%%%%%%%%%%%%%%%%%%%%%%%%%%%%%%%%%%%%%%%%%%%%%%%%%%%%%%%%%%%%%

\end{document}

